\documentclass[a4paper]{article}

\usepackage{baskervald}
\usepackage[utf8]{inputenc}
\usepackage[T1]{fontenc}
\usepackage{textcomp}
\usepackage{amsmath, amssymb}
\edef\restoreparindent{\parindent=\the\parindent\relax}
\DeclareSymbolFont{letters}{OML}{ztmcm}{m}{it}
\usepackage{parskip}
\restoreparindent
\newtheorem{problem}{Problem}
\newtheorem{example}{Example}
\newtheorem{lemma}{Lemma}
\newtheorem{theorem}{Theorem}
\newtheorem{problem}{Problem}
\newtheorem{example}{Example}
\newtheorem{definition}{Definition}
\newtheorem{lemma}{Lemma}
\newtheorem{theorem}{Theorem}
\usepackage{parskip}

\begin{document}

Acabo de volver de tu casa. Mi casa está a oscuras pero, al entrar, siento que
alumbro las tinieblas. Recuerdo cierto pasaje del Evangelio de Juan: «Y una luz
brillaba en las tinieblas, y las tinieblas no la comprendían». 

Hace cierto tiempo, hubiera sido imposible para mí hacerte un lugarcito
en mi corazón. Pienso en aquella chica que me «descubrió», que cometió el
error de confesarme que «veía mi tristeza», y a la que alejé para siempre. Sé
que esto suena ridículo, pero incluso entonces no había nada que yo deseara más
que ser capaz de un amor que fuera verdaderamente puro. Probablemente vos sepas
también que una persona que no puede amar no puede ser feliz---esto es un
hecho---y yo no podía soportar ver que, en un mundo ante el cual yo tanto
deseaba oponer alguna forma de pureza y de ternura, la oscuridad procedía al
menos en parte de mi propio corazón.

Supongo que siento cierto peso particular esta noche porque caigo en la cuenta
de que te conté cosas que no conté a casi nadie. Soy feliz, pero esta
felicidad es una conquista nueva. Tengo dolores que no logro transformar y
permanecen en mí como una quieta forma de melancolía. Mi vida fue tal que, en
cierto punto, sólo una pregunta aparecía ante mí, como si enunciada por una voz
que no era mía: \textit{¿si la vida es sufrimiento, para qué vivir?} Decidí que respondería
a esta pregunta, que lo haría decidida y honestamente, y me juré que si la
respuesta era «\textit{no}» me entregaría a la muerte. (Esta es la única posición
posible.) Ya te he confesado que, en cierto punto, parecía una certeza que debía
morir, y si no fuera por lo que yo, «un chico conectado con la ciencia», sólo
puedo describir como un despertar espiritual, ciertamente estaría muerto hoy.

Sé que mucha gente se reiría de esto. Al mismo tiempo, sé ---estoy seguro--- de
que vos no sos esa gente. Comprendés el sufrimiento, incluso esa rara clase de
sufrimiento que nos conduce al borde de la destrucción.

Supongo que el punto que estoy tratando de hacer es que solía albergar en mí la
absoluta convicción de que no podía amar ni ser amado, no en un sentido
espiritual ---si me permitís ponerlo de ese modo---. Alguna vez viví esa forma
perfecta de amor, pero tuve que enterrar a quien amaba y resignarme a que no la
vería sino en los confines de mis sueños. ¿Quién se atreve a amar, a amar un
amor \textit{puro}, después de eso? No sé, estoy divagando, lo lamento.

Hoy te dije algo que considero cierto: vinimos a este mundo para ser felices, y
quienes son felices hacen la voluntad de Dios. Expresarlo de este modo puede
sonar extraño, pero considero que este es \textit{precisamente} el lenguaje
necesario para hablar de estas cosas. Ser feliz no significa divertirse, como
pensaba esa chica que me impresionó tanto, sino amar a todos, perdonarlo todo,
en fin, seguir ciertos preceptos sencillos que caen dentro de la categoría de
«cristianos». Cuando entendí que la felicidad residía en estas verdades tan
simples, tan extraordinariamente simples que la mente más sencilla podría
entenderlas, y podría deducirlas sin necesidad de abrir un solo libro, mi
corazón se abrió nuevamente y descubrí que tal vez podía atreverme no sólo a
querer sino a ser querido. 

A esta altura del partido te he contado varias cosas acerca de mí, y no
necesariamente las mejores. Has visto ciertas cicatrices, ciertas oscuridades,
incluso tal vez cierta desagradable tendencia a la melancolía. Cada vez que
mostraba esas cosas, para mí se ponía en juego algo vital, y cada vez
respondiste con tenura y compasión. Siento ---no te rías--- que Dios te puso en
mi camino. Si digo esto en un sentido alegórico o literal depende de si me lo
preguntás sobrio o con dos pintas encima: por ahora dejo la frase tal como es,
porque expresa lo que siento.

¿Qué más decirte? Hablé demasiado de mí, es un defecto, soy tristemente
autorreferencial. ¿Sabés algo? Al escucharte cantar, siempre recuerdo ciertos
versos de José Pedroni:

\begin{quote}
    \textit{Si tu alma es como tu voz,}\newline
    \textit{alabado sea Dios!}
\end{quote}

Es probable que vos misma no veas cuánta luz irradiás. Para agotar hasta el
hartazgo las referencias bíblicas, pienso que realmente encarnás aquellas
hermosas palabras: «\textit{ustedes son la luz del mundo}». Yo mismo aspiro a
ser así en la vida y admiro eso de vos. Quisiera aprender de vos.

Son las dos de la mañana, estoy escribiendo hace una hora y mañana curso
temprano. Me retiro sin poder corregir demasiado lo que escribí, posiblemente dé
vergüenza, qué mal gusto tengo...! 

\begin{quote}
Te quiero, te abrazo, te beso. \newline
---$S$
\end{quote}

~ 

PD: Sacate \textit{Gricel} en el teclado.


\end{document}
